% ============================================================================
%  Capitolo 10 -- Introduzione all'elettronica e all'elettronica digitale (Lezione 10)
% ============================================================================
\chapter{Elettronica: di che cosa è fatto un bit}
\label{ch:electronics}

\epigraphline{Qualunque tecnologia sufficientemente avanzata è indistinguibile
dalla magia.}{Arthur C. Clarke, \emph{Profiles of the Future}, 1973}

\begin{chapgoals}
Alla fine di questo capitolo saprete
\begin{itemize}
\item usare tensione, corrente, resistenza e potenza, e applicare la legge di Ohm a un circuito semplice;
\item dimensionare la resistenza di un LED e calcolare un partitore di tensione;
\item spiegare che cos'è un transistor e come un transistor MOS funziona da interruttore;
\item dire come le tensioni diventano valori logici, e che cos'è un margine di rumore;
\item disegnare l'invertitore, il NAND e il NOR in CMOS, e contarne i transistor;
\item spiegare le resistenze di pull-up, gli ingressi flottanti e perché una porta impiega tempo ed energia per commutare.
\end{itemize}
\end{chapgoals}
% Generated by tools/textbook/export_listings.js from tools/decks/mastery.js and tools/concepts/registry.json: do not edit.
\begin{chapstatus}{Stato di apprendimento di questo capitolo}
\masterychip{now}{DA COMPRENDERE OGGI}\enspace gli ingressi con pull-up; tensione, corrente, legge di Ohm; il transistor come interruttore; dai volt ai bit; le porte CMOS.\par\smallskip
\masterychip{use}{DA UTILIZZARE OGGI}\enspace il programma completo, gli strumenti di PRG32 (si usa, non si studia, in questo corso).\par\smallskip
\masterychip{later}{DA COMPRENDERE PIÙ AVANTI}\enspace il clock (Capitolo~\ref{ch:sequential}).
\end{chapstatus}
\begin{chapstatus}{Percorso del capitolo}
Alla fine del capitolo saprete:
\begin{enumerate}[nosep,leftmargin=2em]
\item spiegare come un transistor MOS funziona da interruttore e come CMOS costruisce una porta
\item leggere livelli logici e margini di rumore dalle tensioni
\item dimensionare la resistenza di un LED e spiegare il pull-up di un pulsante
\end{enumerate}
Si appoggia su: solo l'installazione del Capitolo~\ref{ch:setup}. Prossimo passo: Capitolo~\ref{ch:combinational}.
\par\smallskip\textbf{Guardando indietro, guardando avanti.} Nella Lezione 02 il contatore sopravviveva da un frame all'altro in memoria, non in un registro. Oggi una cella di memoria costa circa 6 transistor e un flip-flop circa 20: ricordare un bit è un circuito. La Lezione 13 mostra come latch e flip-flop conservano un bit. \emph{Richiamo:} Dove abitava il contatore della Lezione 02 tra due frame, e perché non in t1?
\end{chapstatus}


Le porte degli ultimi due capitoli erano simboli sulla carta. Questo capitolo
dà loro un corpo fisico, prima che il prossimo le componga in macchine. Non si
presuppone alcuna conoscenza di fisica. I numeri citati per il kit vengono dal
datasheet dell'ESP32-C6 \citep{esp32c6trm} e dalla documentazione hardware di
\prg.

% ----------------------------------------------------------------------------
\section{Riscaldamento sulla macchina}

\asmfile[caption={Un pulsante, un bit.},label=lst:el:button]{code/l10/c03.S}

Fin dai primi capitoli un pulsante premuto è un 1 in un registro
(Listato~\ref{lst:el:button}). Sulla scheda il pulsante è un pezzo di metallo
che unisce un piedino del chip a massa. Premuto significa 0~volt sul piedino,
eppure il programma legge 1. Che cos'è un volt? Perché il piedino rilasciato è
a 3,3~V? E chi decide che 0~V è un bit? Il capitolo segue il pulsante dal dito
al bit, e il suo ultimo paragrafo risponde alle tre domande.

% ----------------------------------------------------------------------------
\section{L'elettricità in quattro grandezze}

\subsection{Tensione, corrente, resistenza, potenza}

La carica elettrica è portata dagli elettroni, e quattro grandezze descrivono
ciò che essa fa in un circuito (Tabella~\ref{tab:el:four}). Un paragone con
l'acqua nei tubi aiuta l'intuizione.

\begin{table}[htb]
\centering\small
\caption{Quattro grandezze.}\label{tab:el:four}
\begin{tabularx}{\textwidth}{@{}llLL@{}}
\toprule
\textbf{Grandezza} & \textbf{Simbolo, unità} & \textbf{Analogia con l'acqua} & \textbf{Sul kit}\\
\midrule
tensione & $V$, volt (V) & la pressione che spinge l'acqua & 3,3 V tra la linea di alimentazione e la massa\\
corrente & $I$, ampere (A) & la portata: litri al secondo & circa 10 mA attraverso un LED acceso\\
resistenza & $R$, ohm ($\Omega$) & un tubo stretto che ostacola il flusso & 45 k$\Omega$: il pull-up dentro ogni piedino di ingresso\\
potenza & $P$, watt (W) & il lavoro che il flusso può compiere ogni secondo & al più 2,5 W da una porta USB 2.0 (5 V, 0,5 A)\\
\bottomrule
\end{tabularx}
\end{table}

Una tensione si misura sempre \emph{tra due punti}. Quando si dice che «il
piedino è a 3,3~V» si intende rispetto alla \term{massa} (GND), il riferimento
comune dell'intero circuito, che è a 0~V per definizione. I prefissi usuali
sono milli (m, $10^{-3}$), micro ($\mu$, $10^{-6}$), nano (n, $10^{-9}$), pico
(p, $10^{-12}$), kilo (k, $10^{3}$) e mega (M, $10^{6}$).

\subsection{La legge di Ohm, e la resistenza di un LED}

\begin{theorem}{Legge di Ohm, 1827}
La corrente attraverso un resistore è proporzionale alla tensione ai suoi capi:
\[ V = R\cdot I . \]
\end{theorem}

Un circuito è un anello chiuso: la corrente esce dalla sorgente, attraversa i
componenti e ritorna. Un LED (Figura~\ref{fig:el:led}) si accende quando è
attraversato da corrente, e quando lo è ai suoi capi c'è una caduta di circa
2~V; senza una resistenza in serie la corrente sarebbe di gran lunga troppo
grande. La resistenza prende ciò che resta, $3,3-2=1,3$~V. Per una corrente di
10~mA,
\[ R=\frac{1,3\ \mathrm{V}}{0,010\ \mathrm{A}}=130\ \Omega, \]
e il valore commerciale più vicino è 150~$\Omega$, che dà
$1,3/150 = 8,7$~mA. La tensione diretta di un LED dipende dal suo colore:
circa 1,8--2,2~V per il rosso, circa 3~V per il blu e il bianco.

\begin{figure}[htb]
\centering
\begin{subfigure}[b]{0.3\textwidth}\centering
\begin{circuitikz}[scale=0.85,transform shape]
  \draw (0,0) node[ground]{} to[battery1,l=3.3<\volt>,invert] (0,3) -- (2.2,3)
    to[R,l=150<\ohm>] (2.2,1.6) to[leDo,l=LED] (2.2,0) -- (0,0);
\end{circuitikz}
\caption{Un LED e la sua resistenza}\label{fig:el:led}
\end{subfigure}\hfill
\begin{subfigure}[b]{0.3\textwidth}\centering
\begin{circuitikz}[scale=0.85,transform shape]
  \draw (0,3.2) node[vcc]{$V_{dd}$} to[R,l=$R_1$] (0,1.6) to[R,l=$R_2$] (0,0) node[ground]{};
  \draw (0,1.6) node[circ]{} -- (1.2,1.6) node[right]{$V_{out}$};
\end{circuitikz}
\caption{Un partitore di tensione}\label{fig:el:div}
\end{subfigure}\hfill
\begin{subfigure}[b]{0.34\textwidth}\centering
\begin{circuitikz}[scale=0.85,transform shape]
  \draw (0,0.2) node[nmos](n){} (n.G) node[left]{G} (n.D) node[above]{D} (n.S) node[below]{S};
  \draw (2.4,0.2) node[pmos](p){} (p.G) node[left]{G} (p.S) node[above]{S} (p.D) node[below]{D};
  \node[lbl] at (0,-1.3) {nMOS};
  \node[lbl] at (2.4,-1.3) {pMOS};
\end{circuitikz}
\caption{I due transistor MOS}\label{fig:el:mos}
\end{subfigure}
\caption{Tre circuiti elementari.}
\end{figure}

\subsection{Serie, parallelo e il partitore di tensione}

Due resistori sono in \term{serie} quando la stessa corrente li attraversa
entrambi: le resistenze si sommano, $R=R_1+R_2$. Sono in \term{parallelo}
quando ai capi di entrambi c'è la stessa tensione: la corrente si divide, e la
resistenza totale, $1/R=1/R_1+1/R_2$, è minore di ciascuna delle due.

Due resistori in serie dividono la tensione (Figura~\ref{fig:el:div}):
\begin{equation}\label{eq:el:div}
V_{out} = V_{dd}\,\frac{R_2}{R_1+R_2}.
\end{equation}
Resistori uguali danno metà dell'alimentazione; una $R_2$ molto più piccola
porta $V_{out}$ vicino alla massa. La formula segue dalla legge di Ohm, poiché
la corrente $I=V_{dd}/(R_1+R_2)$ attraversa $R_2$, e dalle leggi di Kirchhoff:
le correnti che entrano in un nodo sommate danno quelle che ne escono, e le
tensioni lungo una maglia sommate danno zero. Il partitore è la chiave del
circuito del pulsante, dove $R_2$ è un interruttore: una resistenza che è o
infinita o quasi nulla.

\subsection{Potenza ed energia}

La \term{potenza} è tensione per corrente, $P=V\cdot I$, in watt; in un
resistore diventa calore, $P=R\cdot I^{2}$. L'\term{energia} è potenza per
tempo, in joule; una batteria immagazzina energia. Il LED di prima consuma
$3,3\ \mathrm{V}\times 8,7\ \mathrm{mA} = 29$~mW, di cui 11~mW scaldano la
resistenza. Ogni transistor che commuta dentro un processore dissipa un po' di
energia, e con miliardi di transistor il calore è il limite dei chip moderni:
il processore di un portatile dissipa decine di watt in pochi centimetri
quadrati.

\subsection{Il condensatore: perché i segnali impiegano tempo}

Un \term{condensatore} immagazzina carica elettrica come un serbatoio
immagazzina acqua; la sua capacità si misura in farad (F). La sua tensione non
può saltare: il condensatore va riempito, o svuotato, attraverso una
resistenza, e il tempo che ciò richiede è dell'ordine del prodotto
\[ \tau = R\cdot C, \]
la \term{costante di tempo}, trascorsa la quale la tensione ha percorso circa
il 63\% della strada. Un piedino di ingresso del kit ha una capacità di circa
2~pF e sale attraverso il suo pull-up di 45~k$\Omega$ in circa
$45\ \mathrm{k}\Omega\times 2\ \mathrm{pF}=90$~ns (i collegamenti e il pulsante
ne aggiungono altra, quindi il valore reale è maggiore).

Ogni filo e ogni ingresso di un transistor è un piccolo condensatore. È da qui
che viene il ritardo di propagazione di una porta: resistenza per capacità è un
tempo, ed è la ragione per cui in un circuito niente è istantaneo. Un
condensatore carico o scarico è anche un modo di memorizzare un bit: la cella
DRAM del Capitolo~\ref{ch:memory}.

% ----------------------------------------------------------------------------
\section{Semiconduttori}

\subsection{Conduttori, isolanti, semiconduttori}

I conduttori, come il rame, lasciano scorrere liberamente la carica; gli
isolanti, come il vetro, no. I \term{semiconduttori}, il silicio sopra tutti,
stanno in mezzo, e la loro conduzione può essere controllata. Il controllo si
ottiene con il \term{drogaggio}: l'aggiunta di pochi atomi di un altro
elemento. Il silicio di tipo~n ha elettroni in più, che sono portatori di
carica negativi; il silicio di tipo~p ha elettroni mancanti, chiamati
\emph{lacune}, che si comportano come portatori positivi. Dove una regione p
incontra una regione n c'è una \term{giunzione}, il mattone di ogni
dispositivo. Un chip è una fetta di silicio con miliardi di regioni drogate,
collegate da strati di metallo. Il silicio è il secondo elemento più abbondante
nella crosta terrestre, l'elemento della sabbia, e un'intera industria poggia
sul purificarlo e drogarlo con estrema precisione.

\subsection{Il diodo e il LED}

Una sola giunzione fa un \term{diodo}, una valvola a senso unico per la
corrente. In senso diretto, da p a n, conduce una volta che la tensione supera
una soglia, circa 0,7~V per il silicio; in senso inverso blocca. Un diodo a
emissione di luce emette luce quando conduce; il suo colore dipende dal
materiale, e così la sua soglia. I diodi trasformano la corrente alternata in
continua in ogni alimentatore. I LED sono dappertutto in un calcolatore:
indicatori, la retroilluminazione del display del Capitolo~\ref{ch:video},
collegamenti ottici; e anche il laser dei dischi ottici del
Capitolo~\ref{ch:storage} è un diodo.

\subsection{Il transistor: un interruttore senza parti mobili}

Un \term{transistor} è un dispositivo con tre terminali in cui un piccolo
segnale su uno controlla la corrente tra gli altri due. Usato come
amplificatore ha reso portatile la radio; usato come interruttore ha fatto il
calcolatore (Tabella~\ref{tab:el:history}). Valvole termoioniche e relè
facevano lo stesso lavoro nei primi calcolatori del Capitolo~\ref{ch:history},
un milione di volte più grandi e più lenti.

\begin{table}[htb]
\centering\small
\caption{Il transistor in sei date.}\label{tab:el:history}
\begin{tabularx}{\textwidth}{@{}lP{0.34\textwidth}L@{}}
\toprule
\textbf{Anno} & \textbf{Chi} & \textbf{Che cosa}\\
\midrule
1947 & Bardeen, Brattain e Shockley, Bell Labs & il primo transistor: sostituisce la valvola termoionica\\
1958--59 & Jack Kilby; Robert Noyce & il circuito integrato: più transistor su un solo pezzo di semiconduttore\\
1959--60 & Atalla e Kahng, Bell Labs & il transistor MOS: piccolo, semplice da costruire, quasi nessuna corrente di controllo\\
1963 & Wanlass e Sah, Fairchild & il CMOS: coppie di transistor MOS complementari che assorbono potenza solo quando commutano\\
1971 & Intel 4004 & 2\,300 transistor MOS: un processore su un chip\\
oggi & l'intera industria & decine di miliardi di transistor CMOS su un chip\\
\bottomrule
\end{tabularx}
\end{table}

\subsection{Il transistor MOS come interruttore}

Il nome MOS elenca i tre strati del dispositivo: metallo, ossido,
semiconduttore. I suoi tre terminali sono il \term{gate}, che controlla se la
corrente può passare, e il \term{source} e il \term{drain}, tra i quali essa
passa (Figura~\ref{fig:el:mos}). Il gate è isolato: non assorbe praticamente
corrente (al più 50~nA su un piedino dell'ESP32-C6), solo la minuscola carica
della sua capacità.

\begin{definition}{nMOS e pMOS}
Un transistor \term{nMOS} conduce (l'interruttore è chiuso) quando il suo gate
è alto, e blocca quando è basso. Un transistor \term{pMOS} conduce quando il
suo gate è basso, e blocca quando è alto.
\end{definition}

I due tipi sono complementari: ciò che chiude l'uno apre l'altro. Il cerchietto
sul gate del pMOS significa, come sui simboli logici, «attivo quando basso».
Un altro fatto decide come vengono usati: un nMOS lascia passare un buon 0, ed
è quindi usato per collegare a massa; un pMOS lascia passare un buon 1, ed è
usato per collegare all'alimentazione.

% ----------------------------------------------------------------------------
\section{Dai volt ai bit}

\subsection{Livelli logici: quando una tensione è un 1?}

Un circuito digitale usa due \emph{intervalli} di tensione: basso è 0, alto è 1
(Figura~\ref{fig:el:levels}). Un'uscita garantisce livelli puliti: sul kit,
sopra 2,64~V per un 1 e sotto 0,33~V per uno 0. Un ingresso accetta intervalli
più ampi: sopra 2,475~V legge 1, sotto 0,825~V legge 0. (Con $V_{DD}=3,3$~V il
datasheet dà $V_{OH}\ge 0,8\,V_{DD}$, $V_{OL}\le 0,1\,V_{DD}$,
$V_{IH}\ge 0,75\,V_{DD}$, $V_{IL}\le 0,25\,V_{DD}$.)

\begin{definition}{Margine di rumore}
Il \term{margine di rumore} è la differenza tra il livello che un'uscita
garantisce e il livello che un ingresso richiede. Un disturbo più piccolo del
margine di rumore non può cambiare il bit.
\end{definition}

Sul kit i margini sono $2,64-2,475=0,165$~V per il livello alto e
$0,825-0,33=0,495$~V per il livello basso. Tra le due soglie di ingresso il
valore è indefinito: i segnali devono attraversare quella fascia in fretta e
non restarci mai.

\begin{figure}[htb]
\centering
\begin{tikzpicture}[x=1cm,y=1.5cm]
  % output column
  \fill[prgaccent!30] (0,2.64) rectangle (2.2,3.3);
  \fill[prgaccent!30] (0,0) rectangle (2.2,0.33);
  \draw[prgsteel] (0,0) rectangle (2.2,3.3);
  \node[lbl] at (1.1,3.55) {\textbf{uscita}};
  \node[lbl] at (1.1,2.97) {1};
  \node[lbl] at (1.1,0.16) {0};
  % input column
  \fill[prgamber!30] (5,2.475) rectangle (7.2,3.3);
  \fill[prgamber!30] (5,0) rectangle (7.2,0.825);
  \fill[prgred!12] (5,0.825) rectangle (7.2,2.475);
  \draw[prgsteel] (5,0) rectangle (7.2,3.3);
  \node[lbl] at (6.1,3.55) {\textbf{ingresso}};
  \node[lbl] at (6.1,2.9) {legge 1};
  \node[lbl] at (6.1,0.41) {legge 0};
  \node[lbl,text=prgred] at (6.1,1.65) {indefinito};
  \foreach \y/\l in {3.3/{3,3 V},2.64/{2,64 V},0.33/{0,33 V},0/{0 V}} \node[lbl,anchor=east] at (-0.1,\y) {\l};
  \foreach \y/\l in {2.475/{2,475 V},0.825/{0,825 V}} \node[lbl,anchor=west] at (7.3,\y) {\l};
  \draw[dashed,prgink] (2.2,2.64) -- (5,2.64); \draw[dashed,prgink] (2.2,2.475) -- (5,2.475);
  \draw[dashed,prgink] (2.2,0.33) -- (5,0.33); \draw[dashed,prgink] (2.2,0.825) -- (5,0.825);
  \draw[<->,prgviolet,thick] (3.6,2.475) -- (3.6,2.64) node[midway,right,lbl,text=prgviolet]{margine 0,165 V};
  \draw[<->,prgviolet,thick] (3.6,0.33) -- (3.6,0.825) node[midway,right,lbl,text=prgviolet]{margine 0,495 V};
\end{tikzpicture}
\caption{Livelli logici dell'ESP32-C6 con un'alimentazione di 3,3 V.}
\label{fig:el:levels}
\end{figure}

\subsection{Perché digitale?}

In un circuito \emph{analogico} il valore esatto della tensione è
l'informazione. Ogni stadio aggiunge un po' di rumore e di distorsione, che si
accumulano, e le copie peggiorano a ogni generazione. I circuiti analogici
restano indispensabili ai bordi di un sistema: sensori, radio, suono. In un
circuito \emph{digitale} conta solo l'intervallo, alto o basso. Ogni porta
ripristina i livelli, così che il rumore non si accumula; le copie sono
perfette, e gli errori possono essere rilevati e corretti
(Capitolo~\ref{ch:combinational}). E due livelli sono facili da distinguere con
transistor veloci, minuscoli ed economici. È la storia del
Capitolo~\ref{ch:information}, vista dal filo.

\subsection{L'invertitore CMOS: una porta NOT con due transistor}

La porta più semplice è fatta di un pMOS verso l'alimentazione e di un nMOS
verso massa, pilotati dallo stesso ingresso (Figura~\ref{fig:el:inv}).

\begin{figure}[htb]
\centering
\begin{subfigure}[b]{0.28\textwidth}\centering
\begin{circuitikz}[scale=0.8,transform shape]
  \draw (0,2.0) node[pmos](p){} (0,0) node[nmos](n){};
  \draw (p.S) -- ++(0,0.3) node[vcc]{3,3 V};
  \draw (n.S) node[ground]{};
  \draw (p.D) -- (n.D);
  \draw (p.G) -- ++(-0.4,0) coordinate (g1) |- (n.G);
  \draw (g1) ++(0,-1.0) node[circ]{} -- ++(-0.6,0) node[left]{$A$};
  \draw (0,1.0) node[circ]{} -- ++(0.9,0) node[right]{$Y$};
\end{circuitikz}
\caption{NOT: $Y=\nott{A}$}\label{fig:el:inv}
\end{subfigure}\hfill
\begin{subfigure}[b]{0.34\textwidth}\centering
\begin{circuitikz}[scale=0.8,transform shape]
  \draw (0,3.6) node[pmos](p1){} (2,3.6) node[pmos](p2){};
  \draw (1,1.6) node[nmos](n1){} (1,0) node[nmos](n2){};
  \draw (p1.S) -- ++(0,0.25) coordinate (v1) (p2.S) -- ++(0,0.25) coordinate (v2);
  \draw (v1) -- (v2); \draw (1,0|-v1) node[vcc]{3,3 V};
  \draw (p1.D) -- ++(0,-0.25) coordinate (o1) (p2.D) -- ++(0,-0.25) coordinate (o2);
  \draw (o1) -- (o2);
  \draw (n1.D) -- (n1.D|-o1) node[circ]{};
  \draw (o2) -- ++(0.7,0) node[right]{$Y$};
  \draw (n1.S) -- (n2.D); \draw (n2.S) node[ground]{};
  \node[left] at (p1.G) {$A$}; \node[left] at (p2.G) {$B$};
  \node[left] at (n1.G) {$A$}; \node[left] at (n2.G) {$B$};
\end{circuitikz}
\caption{NAND: $Y=\nott{A\cdot B}$}\label{fig:el:nand}
\end{subfigure}\hfill
\begin{subfigure}[b]{0.34\textwidth}\centering
\begin{circuitikz}[scale=0.8,transform shape]
  \draw (1,5.2) node[pmos](p1){} (1,3.6) node[pmos](p2){};
  \draw (0,1.4) node[nmos](n1){} (2,1.4) node[nmos](n2){};
  \draw (p1.S) -- ++(0,0.25) node[vcc]{3,3 V};
  \draw (p1.D) -- (p2.S);
  \draw (n1.D) -- ++(0,0.25) coordinate (o1) (n2.D) -- ++(0,0.25) coordinate (o2);
  \draw (o1) -- (o2);
  \draw (p2.D) -- (p2.D|-o1) node[circ]{};
  \draw (o2) -- ++(0.7,0) node[right]{$Y$};
  \draw (n1.S) -- ++(0,-0.2) coordinate (g1) (n2.S) -- ++(0,-0.2) coordinate (g2);
  \draw (g1) -- (g2); \draw (1,0|-g1) node[ground]{};
  \node[left] at (p1.G) {$A$}; \node[left] at (p2.G) {$B$};
  \node[left] at (n1.G) {$A$}; \node[left] at (n2.G) {$B$};
\end{circuitikz}
\caption{NOR: $Y=\nott{A+B}$}\label{fig:el:nor}
\end{subfigure}
\caption{Tre porte CMOS. Nel NAND e nel NOR i gate dei transistor segnati con
la stessa lettera sono collegati insieme.}
\label{fig:el:cmos}
\end{figure}

\begin{center}\small
\begin{tabular}{@{}cccl@{}}
\toprule
$A$ & \textbf{pMOS} & \textbf{nMOS} & $Y$\\
\midrule
0 & conduce & interdetto & 1: unita a 3,3 V\\
1 & interdetto & conduce & 0: unita a massa\\
\bottomrule
\end{tabular}
\end{center}

Esattamente uno dei due transistor conduce: l'uscita è sempre collegata a una
linea, mai a entrambe. Con l'ingresso fermo, nessuna corrente scorre
dall'alimentazione a massa, così che il CMOS (MOS \emph{complementare}) non
consuma quasi niente a riposo. E l'uscita è un livello nuovo e pulito
qualunque sia la qualità dell'ingresso: la porta \emph{rigenera} il segnale.

\subsection{NAND e NOR in CMOS}

\begin{prgconcept}[Serie è AND, parallelo è OR]
Due interruttori in serie conducono quando il primo \emph{e} il secondo sono
chiusi; due interruttori in parallelo conducono quando il primo \emph{o} il
secondo è chiuso. È l'idea di Nakajima e Shannon.
\end{prgconcept}

Nel NAND (Figura~\ref{fig:el:nand}) due nMOS in serie raggiungono la massa
solo se $A$ e $B$ valgono entrambi 1, e due pMOS in parallelo raggiungono
l'alimentazione se uno dei due vale 0. Il NOR (Figura~\ref{fig:el:nor}) è il
duale: il parallelo diventa serie e la serie diventa parallelo, esattamente
come le leggi di De Morgan scambiano AND e OR. In ogni porta CMOS la rete di
pMOS che porta l'uscita in alto e la rete di nMOS che la porta in basso sono
l'una la duale dell'altra, nel senso del principio di dualità del
Capitolo~\ref{ch:boole1}. Ciascuna delle due porte ha quattro transistor;
verificate sui disegni le quattro combinazioni degli ingressi.

\subsection{Quanto costano le porte, in transistor}

\begin{table}[htb]
\centering\small
\caption{Il costo di alcuni circuiti in CMOS statico convenzionale; i conteggi
variano con lo stile di progetto.}\label{tab:el:cost}
\begin{tabular}{@{}lll@{}}
\toprule
\textbf{Circuito} & \textbf{Transistor} & \textbf{Nota}\\
\midrule
NOT & 2 & l'invertitore CMOS\\
NAND, NOR (2 ingressi) & 4 & le porte naturali del CMOS\\
AND, OR (2 ingressi) & 6 & un NAND o un NOR seguito da un invertitore\\
XOR (2 ingressi) & da 8 a 12 circa & secondo il progetto\\
sommatore completo & circa 28 & il classico sommatore CMOS statico\\
cella SRAM, un bit & 6 & due invertitori ad anello e due transistor di accesso\\
flip-flop D & circa 20 & due latch, master e slave\\
\bottomrule
\end{tabular}
\end{table}

Un AND costa più di un NAND (Tabella~\ref{tab:el:cost}): in CMOS le porte
invertenti vengono per prime. È per questo che il NAND è la porta universale
del Capitolo~\ref{ch:boole2} in pratica oltre che in teoria. Minimizzare una
funzione significa risparmiare transistor, e quindi area, energia e tempo. Il
sommatore che sta dietro l'istruzione \ins{add}, trentadue sommatori completi,
è fatto di circa 900 transistor.

\subsection{Commutare costa tempo ed energia}

L'uscita di una porta pilota gli ingressi delle porte successive, che sono
piccoli condensatori. Per cambiare livello l'uscita deve caricarli o
scaricarli attraverso la resistenza dei suoi transistor: è il \term{ritardo di
propagazione} della porta. E ogni cambio di livello costa energia: la potenza
spesa nella commutazione è circa
\[ P_{dyn} \approx C\cdot V^{2}\cdot f, \]
dove $C$ è la capacità commutata, $V$ la tensione di alimentazione e $f$ la
frequenza di commutazione. Abbassare la tensione di alimentazione è ciò che fa
risparmiare di più, perché compare al quadrato: è la ragione per cui i chip
sono passati da 5~V a 3,3~V e a circa 1~V all'interno. Alzare la frequenza del
clock alza il consumo in proporzione: la velocità si paga in calore. Il ritardo
di propagazione ritorna nel Capitolo~\ref{ch:combinational} come cammino
critico, e nel Capitolo~\ref{ch:sequential} come la ragione del clock.

\subsection{Scale di integrazione}

\begin{table}[htb]
\centering\small
\caption{Scale di integrazione.}\label{tab:el:scales}
\begin{tabularx}{\textwidth}{@{}llL@{}}
\toprule
\textbf{Scala} & \textbf{Transistor per chip} & \textbf{Esempi}\\
\midrule
SSI, piccola & fino a circa 100 & poche porte: la serie 74, come quattro NAND in un contenitore\\
MSI, media & centinaia & decodificatori, multiplexer, contatori, registri: le macchine dei Capitoli da \ref{ch:combinational} a \ref{ch:registers}\\
LSI, grande & migliaia & i primi microprocessori e le prime memorie: Intel 4004\\
VLSI, molto grande & centinaia di migliaia e oltre & i processori degli anni Ottanta e Novanta\\
oggi & decine di miliardi & sistemi su un chip: processore, memoria, radio e periferiche insieme\\
\bottomrule
\end{tabularx}
\end{table}

L'ESP32-C6 del kit è un sistema su un chip in un contenitore di pochi
millimetri (Tabella~\ref{tab:el:scales}). I chip della serie 74 ricompaiono nei
prossimi capitoli come la forma storica di ogni blocco costruttivo.

% ----------------------------------------------------------------------------
\section{Ritorno alla macchina}

\subsection{Pull-up: un ingresso che non resta mai flottante}

Un ingresso collegato a niente è \term{flottante}: raccoglie rumore e legge 0 o
1 a caso. Una \term{resistenza di pull-up} lega il piedino a 3,3~V
(Figura~\ref{fig:el:pullup}). Con il pulsante rilasciato non scorre corrente, e
il piedino legge un 1 pulito. Con il pulsante premuto l'interruttore unisce il
piedino a massa: 0~V, uno 0 pulito, mentre attraverso la resistenza scorre una
corrente di $3,3\ \mathrm{V}/45\ \mathrm{k}\Omega = 73\ \mu\mathrm{A}$, troppo
piccola per contare. Il circuito è il partitore di tensione della
(\ref{eq:el:div}) con un resistore inferiore che è o infinito o nullo.

\begin{figure}[htb]
\centering
\begin{circuitikz}[scale=0.85,transform shape]
  \foreach \x/\sw/\v/\cap in {0/nos/{3,3 V: legge 1}/rilasciato, 6.5/ncs/{0 V: legge 0}/premuto}{
    \draw (\x,3.6) node[vcc]{3,3 V} to[R,l_=45<\kilo\ohm>] (\x,2) coordinate (p\sw);
    \draw (p\sw) to[\sw] (\x,0.2) node[ground]{};
    \draw (p\sw) node[circ]{} -- ++(1.2,0) node[buffer,anchor=in](b\sw){};
    \draw (b\sw.out) node[right]{\small \v};
    \node[lbl] at (\x+0.6,-0.9) {pulsante \cap};
    \node[lbl,anchor=east] at (\x-0.1,2) {piedino};
  }
\end{circuitikz}
\caption{Il pulsante del kit con la sua resistenza di pull-up. Il triangolo è
il buffer di ingresso del piedino.}
\label{fig:el:pullup}
\end{figure}

Il firmware abilita il pull-up interno di ogni piedino dei pulsanti, legge i
piedini e inverte i bit, così che nel programma premuto è 1. Di un segnale il
cui stato attivo è il livello basso si dice che usa una logica \term{attiva
bassa} \eng{active-low}.

\subsection{Dal dito al bit}

Le domande del riscaldamento hanno ora una risposta, e la
Tabella~\ref{tab:el:finger} le mette in ordine. Alcuni dei suoi passi
appartengono a capitoli che devono ancora venire.

\begin{table}[htb]
\centering\small
\caption{Dal dito al bit.}\label{tab:el:finger}
\begin{tabularx}{\textwidth}{@{}cLP{0.34\textwidth}@{}}
\toprule
\textbf{Passo} & \textbf{Che cosa succede} & \textbf{Dove è spiegato}\\
\midrule
1 & il dito chiude l'interruttore tra GPIO18 e massa & questo capitolo: il circuito\\
2 & il piedino scende da 3,3 V a 0 V, in molto meno di un microsecondo & questo capitolo: $R\cdot C$\\
3 & il buffer di ingresso vede una tensione sotto 0,825 V: uno 0 logico & questo capitolo: i livelli logici\\
4 & il livello è campionato in un flip-flop del registro di ingresso & Capitolo \ref{ch:sequential}\\
5 & il firmware legge quel registro con un'istruzione di load & Capitolo \ref{ch:memory}: I/O mappato in memoria\\
6 & inverte il bit e lo impacchetta con gli altri in \reg{a0} & Capitoli \ref{ch:boole1} e \ref{ch:io}\\
7 & il programma lo esamina: \code{andi t3, a0, 0x01} & Capitolo \ref{ch:firstcontact}\\
\bottomrule
\end{tabularx}
\end{table}

I contatti di un interruttore reale \emph{rimbalzano} per qualche millisecondo
quando si chiudono. Leggere i pulsanti una volta per frame, ogni poche decine
di millisecondi, nasconde il rimbalzo (Capitolo~\ref{ch:io}).

\subsection{Il kit come circuito}

Il kit della Figura~\ref{fig:setup:kit} può ora essere letto come un circuito.
Una \term{basetta sperimentale} \eng{breadboard} ha fori uniti a file nella
parte inferiore, così che i componenti si collegano senza saldature. Due lunghi
bus portano l'alimentazione, 3,3~V e massa; un bus separato porta 5~V per gli
amplificatori audio. Ogni parte condivide la stessa massa: senza un riferimento
comune le tensioni di due parti non significano niente l'una per l'altra. I
pulsanti vanno da un piedino a massa; il display e gli amplificatori ricevono
segnali digitali da altri piedini. Un piedino di uscita può accendere un LED
attraverso una resistenza; un piedino di ingresso ascolta soltanto.

\begin{prgwarn}
Un piedino dell'ESP32-C6 può fornire alcune decine di milliampere: a un LED
serve la sua resistenza in serie. Non unite mai direttamente a massa un piedino
pilotato alto.
\end{prgwarn}

% ----------------------------------------------------------------------------
% Generated by tools/textbook/export_listings.js from decks/plan.js and code/l10.S: do not edit.
\section{Il programma di questo capitolo}\label{sec:lab:l10}

I listati di questo capitolo sono parti di un solo programma,
\fname{code/l10.S} nel repository del corso, stampato per intero nel
Listato~\ref{lst:lab:l10}. I test automatici del materiale del corso lo
compilano e lo eseguono sull'esecutore senza interfaccia di \prgqt{}, con
ingressi programmati, e controllano il frame; la
Figura~\ref{fig:lab:l10} è il frame scritto dal test.

\begin{description}[style=nextline,leftmargin=1.5em]
\item[Sullo schermo.] Un grande quadrato, grigio mentre LEFT è rilasciato e giallo mentre è premuto. Sul kit LEFT è un interruttore tra GPIO18 e massa.
\item[Che cosa premere.] Premete e rilasciate LEFT.
\end{description}

\begin{figure}[htb]
\centering
\includegraphics[width=0.5\textwidth]{lab_l10.png}
\caption{\fname{l10.S} su \prgqt{}.}\label{fig:lab:l10}
\end{figure}

\asmfile[caption={Il programma completo del capitolo, \fname{code/l10.S}.},label=lst:lab:l10]{code/l10/program.S}

Per eseguirlo su \prgqt{} i comandi sono quelli del
Listato~\ref{lst:lab:l10:run}; il Paragrafo~\ref{sec:prg:testing} dà gli
stessi passi per QEMU e per la scheda.

\begin{shell}[caption={Compilare \fname{l10.S} ed eseguirlo su \prgqt{}.},label=lst:lab:l10:run]
# 1. start PRG32-QT on this computer (it listens on port 8080)

# 2. in the PRG32 folder, with ESP-IDF active: build the cartridge
python3 -m prg32 cartridge build ~/arc12/code/l10.S --portable \
    --entry-prefix l10 --name l10 --out l10.prg32

# 3. send it to PRG32-QT and start it
python3 -m prg32 esp32c6 upload-and-run l10.prg32 \
    --url http://localhost:8080 --slot cart0

# 4. look, press, compare with the previous slide
\end{shell}


% ----------------------------------------------------------------------------
% Generated by tools/textbook/export_listings.js from tools/concepts/assessments.json: do not edit.
\begin{chapstatus}{Il filo del capitolo}
Un bit è una tensione dentro un intervallo garantito, e una porta è qualche interruttore: il pulsante del kit è ora un circuito che capite. Prossimo passo: le macchine costruite con le porte.
\par\medskip\textbf{Biglietto d'uscita} (da due a quattro minuti, da soli, su carta):
\begin{enumerate}[nosep,leftmargin=2em]
\item In due frasi: come funziona da interruttore un transistor MOS, e perché un invertitore CMOS non assorbe quasi corrente quando il suo ingresso è stabile?
\item Sul kit un ingresso legge 1 sopra 2,475 V e 0 sotto 0,825 V; un'uscita dà un 1 sopra 2,64 V. Un filo aggiunge 0,1 V di rumore verso il basso. Un 1 può essere letto male? Che cosa legge un ingresso a 1,6 V?
\end{enumerate}
\end{chapstatus}


\section{Verificate di aver capito}

\begin{quiz}{Un LED fa cadere 2 V e deve portare 5 mA da un piedino a 3,3 V. Quale resistenza serve?}
\item 26 $\Omega$
\item 260 $\Omega$
\item 400 $\Omega$
\item 660 $\Omega$
\end{quiz}

\begin{quiz}{In un invertitore CMOS con l'ingresso a 0:}
\item entrambi i transistor conducono
\item il pMOS conduce e l'uscita vale 1
\item l'nMOS conduce e l'uscita vale 0
\item nessuno conduce e l'uscita è flottante
\end{quiz}

\begin{quiz}{Quanti transistor richiede un AND a due ingressi in CMOS statico?}
\item 2
\item 4
\item 6
\item 8
\end{quiz}

\begin{quiz}{Perché l'ingresso del pulsante del kit ha bisogno di una resistenza di pull-up?}
\item Per limitare la corrente del LED
\item Perché il piedino legga un 1 definito quando il pulsante è rilasciato
\item Per rendere il pulsante più veloce
\item Per proteggere il chip dai 5 V
\end{quiz}

\section{Esercizi}

\begin{exercise}[prima su carta, poi sul kit]
\begin{enumerate}[nosep]
\item Un LED blu fa cadere 3,0 V: quale resistenza dà 2 mA da 3,3 V? Quale
  potenza dissipa la resistenza?
\item Un partitore con $R_1=10$ k$\Omega$ e $R_2=20$ k$\Omega$ su 3,3 V:
  calcolate $V_{out}$. Un ingresso del kit lo legge come 0, come 1, o è
  indefinito?
\item Disegnate un NAND a tre ingressi in CMOS e contatene i transistor.
\item Verificate il disegno del NOR per le quattro combinazioni degli ingressi,
  segnando quali transistor conducono.
\item Su \prg, eseguite la cartuccia delle lampade e collegate con un filo un
  piedino dei pulsanti a massa: spiegate ciò che vedete, e che cosa
  succederebbe con il pull-up disabilitato.
\end{enumerate}
\expected{Su carta: 150 $\Omega$ e 0,6 mW; $V_{out}=2,2$ V, nella fascia
indefinita; sei transistor, tre pMOS in parallelo e tre nMOS in serie. Sul kit:
la lampada si accende esattamente come se il pulsante fosse premuto.}
\end{exercise}

\section{Glossario}

\begin{glossarytable}
voltage & tensione\\
current & corrente\\
resistor / resistance & resistore / resistenza\\
capacitor & condensatore\\
ground & massa\\
semiconductor, doping & semiconduttore, drogaggio\\
gate, source, drain & gate, source, drain\\
logic level & livello logico\\
noise margin & margine di rumore\\
pull-up resistor & resistenza di pull-up\\
floating input & ingresso flottante\\
propagation delay & ritardo di propagazione\\
integrated circuit & circuito integrato\\
breadboard & basetta sperimentale (breadboard)\\
\end{glossarytable}

\section{Conclusioni e letture}

La tensione spinge, la corrente scorre, la resistenza ostacola: $V=R\cdot I$
descrive la maggior parte di ciò che serve a un progettista digitale. Un
transistor MOS è un interruttore controllato da una tensione, e nMOS e pMOS
sono complementari. Un bit è un intervallo di tensioni; i margini di rumore e
la rigenerazione sono ciò che rende affidabili i circuiti digitali. Le porte
CMOS sono coppie di reti di interruttori, e le porte invertenti NAND e NOR sono
le meno costose. Commutare costa tempo ($R\cdot C$) ed energia
($C\cdot V^{2}\cdot f$): i limiti fisici dietro ogni clock.

\paragraph{Letture.} Charles Petzold, \emph{Code}, i capitoli dalla torcia
elettrica alle porte logiche; \citet{patterson2020}, paragrafi 1.5 e 1.7 e
appendice~A; \fname{docs/hardware/hardware.md} nel repository di \prg.
